\section{Ordered parallelity interfaces with polynomial encoded cost}
\label{int:section}

The source parallelity construction already retains the complete base surfaces
of the selected product handles, their fibre twisting, and a bounded collection
of source-labelled arcs where they meet the complementary core.  Its component
profiles record which ports belong to a base, but that set does not specify the
cyclic order of the ports on its boundary.  In particular, several successive
port rectangles can form one attaching disk.  Counting them as independent
attaching disks changes the topology of the proposed gluing.

This section supplies an exact ordered interface operator.  It reconstructs
the ported boundary circles, distinguishes a common endpoint from a nonempty
original-boundary gap, transports the local fibre charts, and groups the ports
into the actual connected attaching intervals and circles.  All of this uses
binary rank data.  The construction does not walk through the potentially
enormous collection of individual normal polygons.

The relevant predecessor is the source geometry in the repository's
\texttt{parallelity.tex}, together with its independently reconstructed
cut-cell reference~\cite{ProveIt2026}.  The orbit-counting engine is the
Agol--Hass--Thurston algorithm, abbreviated AHT~\cite{AHT2006}.  The results below
separate three points: the already established polynomial weighted AHT
operation, a sharper bound for sparse marked weights, and the new extraction
of ordered topological attachments.  No priority claim for sparse weighted
orbit computation is intended.

\subsection{The encoded marked-boundary problem}
\label{int:model}

\begin{definition}[Marked polygonal interval bundle]
\label{int:definition}
An encoded marked polygonal interval bundle consists of a compact, possibly
disconnected surface $B$, an interval bundle $E\to B$, and a polygon-family
description with the following properties.
\begin{enumerate}
\item Polygon copies are indexed by finite integer intervals whose endpoints
are written in binary.  Each family has an explicitly listed finite sequence
of sides and corners.
\item Glued side families have partial affine rank maps of slope $+1$ or
$-1$, an explicit permutation of their two endpoints, and a separate bit
specifying whether the interval fibre reverses.  These data describe genuine
surface and bundle gluings.
\item Exactly $P$ distinguished free edges are individually named by the
integers $0,\ldots,P-1$.  Each occurs at one specified rank and has a local
direction from endpoint zero to endpoint one.  Their union is denoted $A$.
Every other free edge has the single label ``original boundary.''
\end{enumerate}
The original-boundary label in this definition carries no additional marked
points, curves, slopes, or subdivisions that must be preserved.
\end{definition}

In the normal-cut application, the distinguished edges are the source core
ports.  The source constructor accepts its existing validated finite,
connected, compact, orientable triangulation with one torus boundary and an
admissible integral standard normal vector.  It allows empty, disconnected,
and one-sided cutting surfaces.  Its geometry has at most $66T$ ports, where
$T$ is the number of tetrahedra.  If $D$ is the number of normal disks, its
corner universe has at most $32D$ points and is represented by $O(T)$
families and side records~\cite{ProveIt2026}.

The separate fibre bit is essential.  An order-reversing map of a one-point
rank interval has the same point map as an order-preserving one.  Neither
that map nor the endpoint permutation alone specifies fibre transport.
The new construction retains all three pieces of information.

The output concerns precisely those boundary circles of $B$ that meet $A$.
It does not list every unported boundary circle.  Their number and the base
topology remain available from the existing base-profile construction.

\subsection{Dense weighted transport and the sharper sparse invariant}
\label{int:weights}

We first make the complexity accounting explicit.  This also corrects an
overly cautious reading of later repository passages that express their cost
only in terms of realized weighted partitions.  The preceding repository
section \texttt{suffix\_folds.tex} already gives a polynomial bound for those
partitions.  Classical weighted AHT is likewise polynomial
\cite[Theorem~16]{AHT2006}.  There is no missing general dense-weight theorem
at this stage of the source-cut calculation.

\begin{lemma}[The existing suffix-fold bound]
\label{int:dense}
Let a valid forward AHT trace on $[0,N)$ start with at most $k$ pairings and
$R_0$ canonical constant weight runs.  Weights may be signed integer vectors.
After $\tau$ truncations the maintained suffix-only transport has at most
\[
                         R_0+2\tau
\]
weight runs.  If the trace has $E$ events, the total number of emitted
histogram blocks is at most
\[
                       E(R_0+2E+4k+2).
\]
For additive source intervals $[l_j,h_j)$ with weights $v_j$, every
intermediate coordinate has absolute value at most
\[
                  M=\sum_j(h_j-l_j)\lVert v_j\rVert_\infty.
\]
\end{lemma}

\begin{proof}
This is the suffix-specific argument already established in the
repository~\cite{ProveIt2026}; we recall it to specify exactly which native
operation is being bounded.  Suppose a translation truncation retains
$[0,C)$ and uses period $p$, with $b=C-p$.  Each removed point is sent to
\[
               \rho(x)=b+((x-b)\bmod p)\in[b,C).
\]
The retained prefix weights stay in place.  Let $J_-$ and $J_+$ be the old
interior changes strictly below and strictly above $C$, respectively.
Extend the removed suffix weight by zero.  Its discrete derivative is
supported at $C$, $N$, and $J_+$.  Summing over residues transports these
changes into one period.  Consequently every new interior change lies in
\[
                   J_-\cup\rho(J_+)\cup\{b,\rho(N)\}.
\]
The change at $C$, if there was one, disappears.  Thus at most two new
boundaries replace the vanished suffix boundaries, giving an increase of at
most two runs.

For a reversing truncation, the removed suffix is reflected into one
retained interval.  Its old interior changes are reflected individually;
only the two endpoints of the receiving interval can be additional changes.
The same bound follows.  Deletion, trimming, merger, and transmission leave
the point weights unchanged.  A static contraction retains an ordered
subsequence of the old values and cannot increase the number of runs.
Induction proves the first assertion.

No trace operation increases the number of pairings.  The complement of
their domains and ranges has at most $2k+1$ static gaps.  Intersecting these
gaps with a weight partition having $R$ runs produces at most $R+4k+2$
pieces, a deliberately conservative bound.  Summing over at most $E$
contractions and using $R\le R_0+2E$ gives the displayed output bound.

Finally, every original point contribution moves to exactly one current
representative or one already emitted orbit.  Taking absolute values before
adding the contributions proves the mass bound, including signed
cancellation.  Thus the growth in numerical values costs only logarithmically
many additional bits.  None of the quotient computations iterates over the
number of represented points.
\end{proof}

The useful strengthening for ordered interfaces is independent of the
number of truncations.

\begin{theorem}[Sparse additive orbit transport]
\label{int:sparse}
Supply an interval pseudogroup on $[0,N)$ by $k$ partial integral isometries,
and supply $K$ singleton weight atoms
\[
                        (x_j,w_j),\qquad w_j\in\mathbb Z^d.
\]
Locations may repeat and weights may have either sign.  Their weighted orbit
histogram is computable in time polynomial in
\[
                  k,\quad \log(N+2),\quad K,\quad d,
\]
and the input weight bit lengths.  Along a valid AHT trace, the active weight
function is supported at at most $K$ points and has at most $2K+1$ canonical
runs.  There are at most $K$ distinct nonzero output weight vectors, with at
most one further zero-weight row.

If the atoms have the unique scalar weights $w_j=2^j$, every active nonzero
value encodes a disjoint subset of the $K$ labels.  Every nonzero run is a
singleton, every nonzero histogram row has orbit multiplicity one, and all
mask values have at most $K$ bits.
\end{theorem}

\begin{proof}
Temporarily distinguish the atoms even when their numerical weights are
equal.  Maintain an allocation of the atoms that have not yet been emitted
to current universe points.  The weight at a point is the sum of its
allocated atoms.

Pairing deletion, trimming, periodic merger, and transmission do not move
the allocated atoms.  A translation truncation maps each atom in the deleted
suffix to its unique retained representative modulo the carrier period.
A reflection truncation maps each such atom to its unique reflected
representative.  When several atoms arrive at the same point, their weights
are added.  A static contraction emits all atoms at a deleted point as the
weight of that completed orbit, then reindexes the remaining points
monotonically.  No operation creates a second copy of an atom.

Induction proves that the atom subsets at all live points and all completed
orbits are pairwise disjoint.  Hence at most $K$ live points can have nonzero
weight.  A function supported on $K$ integer points has at most $K$ nonzero
singleton pieces and at most $K+1$ intervening zero pieces; coalescing equal
adjacent values can only improve this count.  Every nonzero final orbit
contains at least one atom, so there are at most $K$ such orbits and at most
$K$ nonzero histogram rows.  Arbitrarily many unmarked or cancelling orbits
share the single zero row.

Each weight coordinate is bounded in absolute value by
$\sum_j |(w_j)_i|$.  Its bit length is therefore at most the maximum input
coordinate bit length plus $\lceil\log_2(K+1)\rceil$ and a constant.
The maintained quotient/remainder and reflection operations process the
bounded run list with binary arithmetic.  A trace of $E$ events consequently
has polynomial weighted replay cost in the displayed parameters and $E$.
For example, a conservative implementation-level bound is
\[
             O\bigl(E(k+K)\log(k+K+2)+K\log(K+2)\bigr)
\]
arithmetic operations, with the vector dimension and actual integer operand
lengths included in the bit cost.  The polynomial classical AHT trace bound
then proves the claimed polynomial algorithm~\cite{AHT2006}.

For the unique scalar labels, sums of disjoint subsets of distinct powers
of two have no binary carries.  Distinct occupied points cannot carry the
same nonzero mask, because that would allocate one label to two points.
This also rules out a nonzero constant run of length greater than one and
rules out two distinct output orbits with the same nonzero mask.  The last
assertions follow.
\end{proof}

Thus the new bound can be recorded together with the existing one as
\[
                 R\le\min\{R_0+2\tau,\,2K+1\}.
\]
Its advantage here is that $K$ is the number of distinguished endpoints,
which is bounded by the source size.  It is not the number of normal disks.

\subsection{The label partitions in weighted selectors are polynomial too}
\label{int:selectors}

The existing complete base construction uses weighted least representatives
to mark individual quotient vertices and boundary components.  Its forward
implementation retains intervals of surviving original labels; its
independent checker restores deleted gaps backwards.  Both admit an explicit
bound that can be combined with Lemma~\ref{int:dense}.

\begin{proposition}[Tagged selector size]
\label{int:tagged}
Consider a complete monotone forward trace with $E$ events, at most $k$
pairings, and $R_0$ initial weight runs.  Let $G$ be the total number of
contracted static-gap intervals.  Then
\[
                 G\le(2k+1)E,
\]
and there are at most $1+G$ live original-label intervals at any time.
The forward tagged selector emits at most
\[
          E\bigl(R_0+2E+1+(2k+1)E+4k+2\bigr)
\]
raw tagged pieces.  Its coalesced output and its independent backward
verification have polynomial size and bit cost in the same parameters.
\end{proposition}

\begin{proof}
Initially the live original labels form one interval, or none when the
universe is empty.  Removing a contiguous interval of current ranks can
split one live original-label interval into two.  If instead it crosses
several live intervals, it deletes middle intervals and cuts only the two
end intervals, with no larger increase in their number.  Thus one removed
gap increases the live interval count by at most one.  Retaining a prefix
cannot increase it.  This proves the $1+G$ bound.  Each contraction has at
most $2k+1$ gaps, proving the bound on $G$.

At a contraction the emission routine intersects three ordered partitions:
the constant weight runs, the live source-label intervals in current-rank
coordinates, and the static gaps.  If their sizes are $R$, $L$, and $r$,
the common refinement contains at most $R+L+2r$ relevant pieces.  Substitute
\[
                R\le R_0+2E,\qquad
                L\le1+(2k+1)E,\qquad r\le2k+1
\]
and sum over at most $E$ contractions.  This gives the stated conservative
bound.  Sorting and coalescing these explicitly bounded pieces is
polynomial, and their endpoints remain binary source labels.

For the independent reverse construction, the selected tagged intervals
remain disjoint.  Inserting a removed gap has one insertion anchor, which
can split at most one currently selected interval.  Across all reverse
insertions there are at most $G$ such new splits beyond the pieces originally
recorded at contractions.  Truncations insert no selected points: a deleted
suffix contains no least original representative.  Binary gap insertion,
sorting, and coalescing therefore also have polynomial cost.  This argument
applies directly to the producer's live-label operations and the verifier's
separate backward reconstruction~\cite{ProveIt2026}.
\end{proof}

\begin{corollary}[Polynomial source-parallelity accounting]
\label{int:source-polynomial}
For the validated normal-cut source class with $T$ tetrahedra and $D$ normal
disks, the source-parallelity base construction, its tagged selectors, and
the ordered interface construction below have a polynomial-time realization
in $T$ and $\log(D+2)$.  No extra hypothesis about an uncontrolled number
of weighted or label partitions is required.
\end{corollary}

\begin{proof}
The source geometry has $O(T)$ records and binary rank endpoints.  Its
corner and polygon universes have sizes bounded by constant multiples of
$D$~\cite{ProveIt2026}.  The first two source-parallelity queries are the
weighted corner and boundary selectors.  Proposition~\ref{int:tagged}
bounds their outputs polynomially in their AHT traces.  The subsequent
mark-placement operation intersects their ordered representative intervals
with the $O(T)$ polygon-carrier intervals.  A common refinement of two
ordered interval families has polynomial size.  Hence the weights supplied
to the base-component and fibre-cover histograms have polynomially many
input intervals.  Lemma~\ref{int:dense} bounds both histograms and their
integer sizes.  Finally, Theorem~\ref{int:sparse} applies to the two ordered
interface queries.  Choosing classical polynomial AHT schedules for these
finitely many queries proves the corollary.
\end{proof}

The new ordered-interface implementation explicitly selects the original
AHT periodic-overlap threshold with \texttt{periodic\_rule='aht'} in both
queries.  The pre-existing source-parallelity implementation uses its
documented forward AHT engine, including its already justified stronger
Fine--Wilf merger.  The partition bounds above apply to either valid forward
schedule; their conversion to a worst-case input bound invokes the
polynomial trace bound for the chosen schedule.  In particular, the existing
phrases ``polynomial in the realized partitions'' need not be interpreted as
an additional unresolved source-parallelity hypothesis.  This is a
clarification of local encoded operations.  It gives no bound on the number
of cutting surfaces or hierarchy states that a complete recognizer must
discover.

\subsection{Two marked orbit queries recover the boundary words}
\label{int:construction}

Give every individual polygon corner its place in a single integer universe
$[0,N)$.  A gluing of polygon sides identifies their endpoint corners.  The
source rank maps and corner offsets express those identifications as partial
integer isometries.  Let $V$ denote the resulting interval equivalence
system.  Its orbits are the actual quotient vertices of the polygonal
surface.

For port $p$, introduce the two endpoint tokens
\[
                         e=2p,\qquad e=2p+1.
\]
Tokens are distinct names even when their corner points are identified or
are literally the same source corner.  Place the singleton scalar weight
$2^e$ at the corresponding corner.  The first query is the weighted orbit
histogram of $V$.  It records which distinguished endpoints are the same
surface vertex.  At a boundary vertex at most two port endpoints can occur,
because the boundary is a one-manifold.  Accordingly, each nonzero mask has
one or two bits and occurs once.

Next add to $V$ the endpoint identifications across every original-boundary
free edge, omitting every core-port edge.  Call the resulting system $G$.
It records connectivity along the boundary after the interiors of the port
edges have been removed.  Its marked components are intervals, including
the possible one-vertex interval between consecutive port edges.  Each has
exactly two endpoint tokens.

We do not need to query $G$ separately.  Instead, construct its two-sheeted
fibre cover $\widetilde G$ on $[0,2N)$.  Across an original corner gluing,
the source fibre bit either preserves or exchanges the sheets.  Across an
original-boundary free edge, the local polygon fibre chart preserves them.
The lifted token $(e,s)$ receives weight
\[
                              2^{2e+s},\qquad s\in\{0,1\}.
\]
The second query is the weighted histogram of $\widetilde G$.  Every marked
interval has two disjoint lifts, so the nonzero masks give a perfect
matching of the $4P$ lifted tokens.  If
\[
                            (e,0)\longleftrightarrow(f,b),
\]
the matching identifies the gap endpoints $e,f$ and its fibre reversal $b$.
The other lift must be
\[
                            (e,1)\longleftrightarrow(f,1-b).
\]
Both relations are checked; in particular $e\ne f$.

\begin{theorem}[Exact ordered interface extraction]
\label{int:ordered}
For an encoded marked polygonal interval bundle as in
Definition~\ref{int:definition}, two sparse weighted orbit queries determine,
for every boundary circle meeting $A$:
\begin{enumerate}
\item the cyclic order and local directions of all its port edges;
\item whether each intervening original-boundary gap is nonempty or consists
of a common vertex;
\item the fibre reversal along each gap in the specified endpoint charts;
\item the connected components of $A$, classified as intervals or whole
boundary circles, with the endpoint charts and fibre transport of every
interval component.
\end{enumerate}
The output has $O(P)$ records.  The queries have $2P$ and $4P$ singleton
tokens, respectively, and their weight partitions have at most $4P+1$ and
$8P+1$ runs.  The construction has polynomial encoded cost.  If the total
interval bundle is orientable, the fibre parity around every output
boundary circle is zero.
\end{theorem}

\begin{proof}
The first query computes actual corner equivalence and therefore tells
exactly which endpoint tokens are the same surface vertex.  Distinct
token names give disjoint bits, so their masks lose no endpoint information.
Unmarked vertices contribute only zero weight, regardless of their number.

Consider a boundary circle that contains at least one port.  Removing the
interiors of its port edges leaves a finite collection of compact intervals.
An interval can be a single vertex when two port edges are successive with
no intervening original-boundary edge.  Its two boundary incidences are the
two endpoint tokens that it carries.  The system $G$ gives exactly these
interval components: its corner relations identify surface vertices, and
its additional relations cross precisely the non-port boundary edges.
Unported boundary circles and interior vertices carry no endpoint labels.

An interval bundle over an interval is trivial.  In the two-sheet fibre
cover its two endpoint lifts therefore form two disjoint pairs.  The second
histogram recovers these pairs and the relative fibre chart.  This proves
the asserted endpoint matching and fibre reversal.  A gap is a single
vertex if and only if its two tokens belong to the same $V$-orbit.  Indeed,
the forward implication is immediate.  For the converse, a positive
original-boundary path beginning and ending at that vertex, in addition to
the two incident port-edge ends, would give more than two boundary
incidences there.  That contradicts the boundary one-manifold condition.

Join the recovered gap matching to the fixed port matching
$2p\leftrightarrow2p+1$.  Each endpoint is incident to one gap and one port,
so these alternating components are exactly the ported boundary circles.
Start each component at its least port, entering that port through its
endpoint zero.  Crossing a port in direction zero to one writes $p+1$;
crossing it in the other direction writes $-(p+1)$.  Crossing the following
gap determines the next port endpoint.  This produces a deterministic
signed cyclic word, up to the normalization just specified.  Every port
occurs exactly once in one such word, and there are exactly $P$ gaps
between consecutive ports in the aggregate.

Split a cyclic word at every nonempty original-boundary gap.  Its remaining
maximal runs are precisely the connected interval components of $A$.
If there is no such gap, the entire word is one circle component of $A$.
The local fibre transport along a port is trivial in its own polygon
chart.  Consequently, the XOR of the intervening gap bits gives the fibre
transport of each interval run, and the XOR around a cyclic word gives
its fibre monodromy.  When the total space $E$ is orientable, its interval
orientation character is $w_1(B)$.  The restriction of $w_1(B)$ to every
boundary circle vanishes, since that circle has an annular collar in $B$.
The total XOR around each boundary word is therefore zero.

Both orbit systems have a number of pairing families bounded by a constant
multiple of the input side and gluing records.  Doubling sheets doubles
the universe and only multiplies the number of families by a constant; it
does not expand the universe.  Theorem~\ref{int:sparse}, applied to $2P$
and $4P$ singleton tokens, gives the stated run and complexity bounds.
Decoding masks, forming the two matchings, traversing the $P$ port edges,
and splitting the words use polynomial work and $O(P)$ output records.
\end{proof}

\begin{corollary}[Source-size bound for ordered normal-cut interfaces]
\label{int:normal}
For the source normal-cut geometry, the reduced ordered core interfaces have
$O(T)$ records and can be computed in time polynomial in $T$ and
$\log(D+2)$, independently of the numerical normal multiplicities.
\end{corollary}

\begin{proof}
Substitute the source bounds $P\le66T$, $N\le32D$, and $O(T)$ polygon-side
records into Theorem~\ref{int:ordered}.  The geometry and the source-normal
validation are binary rank operations.  If $D=0$, there are no ports and the
empty interface follows directly from the source construction.
\end{proof}

\subsection{Attaching intervals, attaching annuli, and the exact invariant}
\label{int:attachments}

\begin{corollary}[Connected attachment types]
\label{int:attachment-types}
Suppose the output contains $a$ interval components and $c$ circle components
of $A$.  Then
\[
                \chi(A)=a,\qquad
                H_1(A;\mathbb Z)\cong\mathbb Z^c.
\]
In the orientable total-bundle setting, the interval-bundle attachment above
each component is a disk for an interval component and an annulus for a
circle component.  Adjacent source ports in one interval run form one
attaching disk.
\end{corollary}

\begin{proof}
A finite disjoint union of compact intervals and circles has the stated
Euler characteristic and first homology.  The restriction of any interval
bundle to an interval is a product, hence a disk.  On a boundary circle of
an orientable total bundle the fibre character is zero by
Theorem~\ref{int:ordered}; its vertical restriction is therefore an annulus.
The decomposition into components is the maximal-run decomposition already
proved, so the several rectangles over consecutive port arcs constitute
one connected disk, rather than separate attaching disks.
\end{proof}

The value $a$ agrees with the existing independently checked attachment
Euler characteristic, computed there as the number of quotient vertices
incident to core ports minus the number of core ports.  The new operator
retains the extra information that this Euler characteristic omits:
the number and identities of entire attaching circles, cyclic order,
individual interval membership, and fibre transport.

For example, the four-tetrahedron layered-meridian fixture has $15$ ports
but only three interval attachments.  An especially small annular
obstruction uses the two-tetrahedron layered source and the coordinate rows
\[
                (0,0,1,1,1,0,0),\qquad (1,1,1,1,0,0,0).
\]
Its source parallelity base is the edge-cap-only disk already described in
the repository's parallelity analysis~\cite{ProveIt2026}.  The new output
has the single word
\[
                         (1,2,3,4,5,6,7,8)
\]
and every following gap is a common vertex.  Thus $a=0$, $c=1$: its whole
base boundary is attached, and the corresponding vertical attachment is an
annulus.  A criterion that permits only attaching disks must exclude this
case.  It cannot replace the circle by eight independent disk attachments.

\subsection{Replication along a two-sided normal surface}
\label{int:replication}

The dependence on a large common multiplicity can be separated even more
strongly when a global coorientation is available.  The following statement
concerns the source decomposition itself, including its marked attachment
data; it is not merely a statement about the number of complementary
components.

\begin{theorem}[Positive-multiplicity replication]
\label{int:replication-theorem}
Let $F$ be a nonempty two-sided properly embedded normal surface with
admissible standard coordinate vector $x$ in a triangulated orientable
$3$-manifold satisfying the source contract.  Fix a global coorientation
on each component of $F$.  Write $K_x$ for the source complementary core
and $E_x\to B_x$ for its selected parallelity bundle.  For every integer
$m\ge1$, the standard normal surface with coordinates $mx$ admits the
following source-compatible description.
\begin{enumerate}
\item Its core is canonically identified with $K_x$ by the local handle
types and owners.  Source rank labels on its interface are obtained by
multiplying the corresponding local gap ranks by $m$.
\item In every product-handle family, whose local gap ranks have the form
$1\le r<md$ after replication, retain exactly the copies with $m\mid r$.
These copies form a union of complete bundle components.  Division of
their local ranks by $m$ identifies their polygonal bundle, including its
fibre charts and marked core attachment, with $E_x\to B_x$.
\item Every source core port belongs to this divisible-rank subbundle.
The remaining product handles form $m-1$ unported copies of $F\times I$,
with $m-1$ base copies of $F$.  If $F$ has $q$ connected components, this
means $(m-1)q$ additional connected unported bases.  Existing unported
components of $B_x$ are retained as well.
\item The reduced signed core-port words, the zero-or-positive
original-boundary gap indicators, and the fibre transport are unchanged
under the source identification.  In particular, the numbers and owner
incidences of attaching intervals and attaching circles are unchanged.
\end{enumerate}
The identification is by source family and local rank; the flattened
integer offsets and the numerical polygon-family IDs need not agree.
Topologically,
\[
 \operatorname{Cut}(M,mF)\cong
 \operatorname{Cut}(M,F)\;\sqcup\;
 \bigsqcup_{j=1}^{m-1}(F\times I).
\]
\end{theorem}

\begin{proof}
Choose a small two-sided normal collar of $F$ compatible with the fixed
coorientation.  Replace each local disk of $F$ by $m$ parallel disks in
that collar.  Face matching follows the old matching and the collar
coordinate, so this is an embedded normal representative of $mx$.
There are $m-1$ intervening product slabs over every component of $F$.
The remaining part of the complement is the old cut complement, with
its two cut sides shifted to the two exterior copies of the collar.
This proves the displayed topological decomposition.  We now verify
its stronger assertion about the specified source handles.

A local product handle lies between consecutive objects in one ordered
stack: parallel normal disks in a tetrahedron, parallel normal arcs in
a face, or normal intersection points on an interior edge.  If the old
stack has size $d$, its product ranks are $1,\ldots,d-1$.  Replication
replaces each old object by a block of $m$ objects.  The gap between
successive old blocks has rank $mr$, whereas an interior gap within a
block has rank not divisible by $m$.  Thus the divisible ranks are
exactly the old inter-block product handles.  If an old stack has
$d=1$, its source family was absent; the new ranks $1,\ldots,m-1$
are all nondivisible, as required.

This distinction is respected by every source gluing.  Expressed in
\emph{local gap ranks}, each gluing is
$r\mapsto\varepsilon r+c(x)$, where $\varepsilon\in\{1,-1\}$ and
$c(x)$ is a homogeneous integral linear combination of the relevant
normal disk counts.  For example, the prism-to-face maps are $r$,
$t_v+r$, or $t_v+q-r$, where $t_v$ is a triangle count and $q$ the
relevant quadrilateral count.  The prism-to-edge and face-to-edge maps
have the forms $t_v+r$, $w-t_v-r$, $r$, or $w-r$, where $w$ is an
edge weight; restrictions choose the appropriate family and interval.
These are gap ranks, so the $-1$ that occurs in a zero-based
\emph{disk} reflection has already been removed by the source conversion
to gap coordinates.  After replacing $x$ by $mx$, the offset becomes
$mc(x)$.  Hence divisible ranks glue only to divisible ranks, and
division by $m$ recovers every original gluing and its domain.  The
endpoint permutations and fibre-reversal bits depend on source
orientations and local disk types, which have not changed.  This proves
the asserted bundle identification, including local fibre charts.

The complementary source $0$-handles are the chambers at stack ends
and at changes of disk type.  Its $1$-handles are the corresponding
nonparallel face regions, its $2$-handles are the two exterior edge
gaps (or the sole gap for zero edge weight), and its $3$-handles are
the fixed interior vertex regions.  Positive scaling preserves all
zero-or-positive tests and disk-type supports.  Consequently their
types, owners, and attaching charts identify with those for $x$.
Each core-facing product port occurs at a transition rank such as
$t_v$, $t_v+q$, an arc count, or its complement in an edge weight.
These ranks are homogeneous in $x$, so all new port ranks are
divisible by $m$ and all old ports reappear at their rescaled ranks.
The possibly newly present product families with old stack size one
have no such interior divisible rank and contribute no core port.

The product handles at the nondivisible ranks fill precisely the
slabs inside the replicated collars.  The chosen global coorientation
orders the $m$ sheets consistently on each component of $F$, thereby
identifying the slabs as $F\times I$ with indices
$1,\ldots,m-1$.  A reversed local stack can exchange residue $s$
with $m-s$; this local effect does not exchange the globally indexed
slabs after the coorientation is used.  Their bases are the corresponding
copies of $F$ and their only vertical boundary lies over $\partial F$
in the original boundary of $M$.  None touches the source core.

Finally, the divisible-rank identification preserves all free-edge
labels, corner identifications, and fibre bits.  It therefore preserves
the two pairing problems used by Theorem~\ref{int:ordered} on every
ported component.  The reduced boundary words, their gap types, and
their attaching-component decomposition are identical after applying
the source identification.  This proves all four assertions.
\end{proof}

\begin{corollary}[Replication of a certified disk cut]
\label{int:replicated-disk}
If $F$ is a connected two-sided normal disk and its primitive source cut
has a valid handlebody census, then every positive multiple has the same
ported core and contributes $m-1$ additional unported balls.  In particular,
if the primitive cut is a single ball, the cut along $mF$ consists of
exactly $m$ balls.
\end{corollary}

\begin{proof}
Apply Theorem~\ref{int:replication-theorem} and use
$D^2\times I\cong D^3$ for the additional unported slabs.
\end{proof}

This theorem permits a future certificate format to retain a primitive
interface proof and a binary replication factor, once a source-bound
coorientation certificate has been checked.  The delivered implementation
still runs its general compressed queries on the supplied vector; it does
not use such a replication shortcut.  The coorientation hypothesis is
substantive: for a one-sided component the local sheet order can reverse
around a loop, and the asserted list of $m-1$ copies of $F$ does not follow.
The general ordered-interface operator itself remains valid in that case.

For the layered-meridian experiments below, the theorem proves independence
of the positive scalar $m$ once the primitive disk and its source interface
have been certified.  It does not prove an interface formula for all
numbers of layers.  In particular, the observed formulas for the number
of primitive ports and attaching intervals as the layer count varies
remain finite experimental observations here; no induction establishing
an all-layer primitive-interface theorem is claimed.


\subsection{Source-bound certificates and implementation contract}
\label{int:certificates}

The producer is
\texttt{normal\_parallelity\_boundary\_words}; the independently implemented
checker is \texttt{verify\_parallelity\_boundary\_words}.  The result retains
the source port records, including their core-handle owners, and contains
the following finite data.
\begin{itemize}
\item Each gap specifies its two endpoint token IDs, whether original
boundary is present, and its fibre reversal.  Endpoints are ordered by ID.
\item Each ported boundary circle contains its signed word, the corresponding
list of following gap IDs, and the whole-core-circle indicator.
\item Each connected attachment contains its interval-or-circle type, its
signed subword, and, for an interval, its entry and exit tokens and fibre
transport.  A least port ID serves as an anchor for joining it to the
existing base profile whose port mask contains that ID.
\end{itemize}

All ports in one connected attachment belong to the same base component.
A base profile containing a port has multiplicity one in the existing
source classification.  Thus the anchor joins the ordered interface to
the base topology without making any choice among repeated unported
components.

The certificate has schema
\texttt{parallelity-boundary-words-v1}.  It binds the input fingerprint,
the proposed interface, the weighted vertex-partition proof, and the
weighted gap-cover proof.  Verification reconstructs the parallelity
geometry through the separate raw-cut-derived reference.  It independently
forms the two pairing systems, checks the supplied AHT traces and weights,
and reconstructs the proposed words and connected attachments.  It does not
import the producer's query builder, mask decoder, traversal routine, or
discovery scheduler.

When $P=0$, the output is the empty interface.  Both proof fields are
explicitly null; neither orbit query is run.  The checker reconstructs the
source and verifies the absence of ports, so this shortcut does not rely
on a caller assertion.  It says nothing about the number of unported base
boundary circles, which this API has not been asked to enumerate.

A discovery-cycle limit is shared by the two queries.  Exhaustion returns
an inconclusive status without an interface or certificate.  A verification
operation limit counts the two supplied ordinary traces.  Cooperative
cancellation propagates through both paths.  The operation returns geometric
information about the supplied source and cutting vector; it does not
itself issue a knot verdict.

\subsection{Independent finite checks and binary scaling}
\label{int:evidence}

The added test module passed all eight test methods in the recorded run.
It covers $193$ source/vector cases: all $157$ saved native source cases,
$16$ layered-meridian cases, eight one-sided cases, and twelve source
relabellings.  Of these, $182$ are compared with a literal expanded oracle.
That oracle expands the actual free polygon edges, identifies corner fibre
charts with a parity union-find, and walks each old-boundary edge one at a
time.  It separately computes connected components using only the core
edges, rather than reusing the producer's maximal-word-run routine.  Eleven
very large native cases use compressed independent source replay and the
attachment-Euler and base-profile checks instead of exceeding the literal
oracle's point cap.

Further tests reject altered source fingerprints, core owners, signed
words, zero-gap flags, fibre bits, endpoint masks, and attachment types.
They verify exact hexadecimal JSON replay with multiplicity $2^{16384}$,
shared caps, cancellation, and producer-disabled verification.  Four
generic sparse-token examples compare weighted translation orbits with
their exact residue classes, up to a $16384$-bit universe.  The test log is
\texttt{boundary-words-tests.txt}; the recorded run took $69.940$ seconds.

The separate benchmark uses actual layered-solid-torus normal coordinates
with $T\in\{1,2,4,8\}$ and
\[
                    m=2^e,\qquad
                    e\in\{0,16,64,256,1024,4096,16384\}.
\]
All $28$ certificates are independently replayed.  Selected measurements
are shown in Table~\ref{int:timings}.  They measure the new ordered-interface
operator, not a complete unknot-recognition run.  They are individual scale
measurements, not statistical or competitive timing claims.

\begin{table}[htbp]
\centering
\small
\begin{tabular}{rrrrrrr}
\hline
$T$ & $m$ & $P$ & $a$ & Produce (s) & Verify (s) & Certificate bytes \\
\hline
2 & $1$ & 2 & 2 & 0.00342 & 0.00371 & 5,949 \\
2 & $2^{16384}$ & 2 & 2 & 0.17363 & 0.09245 & 4,541,499 \\
4 & $1$ & 15 & 3 & 0.10266 & 0.08041 & 47,030 \\
4 & $2^{16384}$ & 15 & 3 & 1.12275 & 0.65597 & 12,861,283 \\
8 & $1$ & 47 & 3 & 0.70518 & 0.53165 & 164,528 \\
8 & $2^{16384}$ & 47 & 3 & 4.47206 & 2.50908 & 30,982,516 \\
\hline
\end{tabular}
\caption{Exact source-bound ordered interfaces.  Here $a$ is the number of
connected interval attachments, not the number $P$ of source ports.}
\label{int:timings}
\end{table}

At $T=2$ and $m=2^{16384}$, the largest vertex and cover weight partitions
have $9$ and $17$ runs, meeting the respective sparse upper bounds.  At
$T=4$ they have $49$ and $97$ runs, and at $T=8$ they have $145$ and $289$.
These remain below $4P+1$ and $8P+1$.  The one-tetrahedron examples have
no ports and use the no-query shortcut; their boundary-interface certificate
has $294$ bytes at every measured exponent.

The distinction between output structure and proof size is visible in the
table.  The number of words, ports, and connected attachments stays fixed
for each source family.  Binary endpoints in the local-reduction traces
still grow with the multiplicity bit length.  In particular, the
eight-tetrahedron extreme example has a roughly $31$-million-byte
certificate, despite its $47$ ports.  The data support polynomial encoded
scaling and constant combinatorial interface size in these families; they
do not support a claim of constant certificate size.

\subsection{An alternative certificate using only three orbit counts}
\label{int:three-count}

The following elementary result supplies an alternative way to authenticate
the marked partitions.  It is useful both as a separate proof principle and
as a possible future checker.  It is not an implemented backend of the
present package.

\begin{theorem}[Three-count certificate for a marked partition]
\label{int:partition-certificate}
Let an interval equivalence system have $C$ orbits and let $K>0$ marked
tokens be placed at specified points.  A proposed partition of the tokens
into $m$ nonempty blocks is exactly their original orbit partition if and
only if the following two tests hold.
\begin{enumerate}
\item Add singleton pairings forming a spanning tree within each proposed
block.  The resulting orbit count is still $C$.
\item In a separate copy of the original system, choose one representative
token from each proposed block and add singleton pairings forming a tree
on these $m$ representatives.  The resulting orbit count is $C-m+1$.
\end{enumerate}
Thus the original count and these two augmented counts form a three-trace
certificate of the marked partition.  The $K=0$ case is immediate and
requires no partition query.
\end{theorem}

\begin{proof}
Adding a pairing can only merge original orbits; it cannot split one.
Therefore the first augmented count equals $C$ precisely when every edge
of every within-block tree has its ends in the same original orbit.  By
connectivity of a tree, all tokens in each proposed block then belong to
one original orbit.

Let $q$ be the number of distinct original orbits occupied by the chosen
representatives.  Connecting all representatives into one tree merges
exactly those $q$ orbits into one, irrespective of repetitions among the
representatives' original orbits.  The second augmented count is therefore
$C-q+1$.  Requiring it to be $C-m+1$ is equivalent to $q=m$, so the proposed
blocks belong to distinct original orbits.  Together with the first test,
this proves exactness.  Conversely, an exact partition plainly satisfies
both count identities.
\end{proof}

There are at most $K-1$ added singleton pairings in either augmentation.
For discovery, equality of two marked points can likewise be tested by
adding one pairing and checking whether the orbit count drops.  Performing
at most $\binom K2$ such tests constructs the entire partition in polynomial
time from the unweighted AHT oracle.  For sparse integer weights, summing
the weights over these marked blocks then recovers their histogram; the
remaining unmarked orbits, together with any zero-sum marked blocks, supply
the zero row.  This gives a slower but conceptually independent polynomial
alternative to Theorem~\ref{int:sparse}.  The shipped two-query construction
uses the more economical sparse weighted transport and its separate local
replay checker.

\subsection{What the ordered interface does and does not settle}
\label{int:limits}

Replacing each nonempty unmarked boundary path by one recorded gap preserves
the marked boundary pair and the relative fibre charts needed for the
interval-attachment calculation.  If a later hierarchy has additional
pattern marks inside those paths, those marks must be distinguished before
compression, with their encoded number included in the input bound.  The
present one-color original-boundary contract cannot silently preserve such
extra data.

The construction also does not classify arbitrary annular attachments,
make the inherited core admissible, produce a new compressing disk, or bound
the number of hierarchy restarts.  Its concrete advance is a source-bound,
polynomial extraction of the exact connected and ordered attachments that
the next topological gluing operation requires.  Together with the dense
and tagged-selector bounds above, it establishes local encoded cost without
turning that local statement into a global quasi-polynomial recognition
claim.
